% Chapter 1, typeset from lectures/L01: the README is the opener and the closing review, and each
% appendix is a section, in the same order and with the same subsections.

\chapter{Circuits, units, and the nodal solver}
\label{c1:ch:circuits}

{\large\itshape Circuit analysis as something a program does, built before there is anything
interesting to analyse.\par}

\begin{chapterbox}{In this chapter}
\begin{itemize}
  \item Charge, current, voltage and power, and the sign convention every later result depends on.
  \item Series and parallel, and the voltage divider as the one circuit worth memorising.
  \item Kirchhoff's two laws, which are the only physical content in the whole of circuit theory.
  \item Thevenin and Norton, and loading as the reason either of them matters.
  \item Nodal analysis: one unknown per node, one equation per node, and the conductance stamp.
  \item Why a program that solves circuits is easier to write than the equations are to solve by
    hand.
  \item Building \code{ael::net} and \code{ael::mna}.
\end{itemize}
\end{chapterbox}

This chapter is the circuit theory the rest of the book needs, and nothing more than that. Most of
it will be familiar if you have an electronics background, and two parts of it will not be used the
way you expect. The first is the \textbf{sign convention}: current into a component is positive,
and every gain, every bias point and every output resistance in Part~2 inherits that choice, so it
is fixed here rather than argued about in Chapter~\ref{c7:ch:smallsignal}. The second is
\textbf{loading}. A divider gives 1.71 V until something is connected to it, at which point it
gives 1.09 V, and that one subtraction is the single arithmetic fact this book returns to six
times, in six different costumes, ending with an operational amplifier that loses 11 dB to it.

Then you write a solver. \Secref{c1:app:b} has you implement nodal analysis before you have met a
single active device, which costs an hour and pays for the rest of the book: from then on every
claim the text makes about a circuit can be checked against a program you wrote, and no result has
to be taken on authority. The method itself is one idea --- one unknown per node, one current-law
equation per node, and a \emph{stamp} that turns each element into matrix entries without anyone
deriving anything.

Work the chapter in order. \Secref{c1:app:a} is the theory, and if you have the background it is a
twenty-minute skim, except for \secref{c1:sec:a6}. \Secref{c1:app:b} builds the solver, and
\secref{c1:sec:b5} is its specification. The exercises in \secref{c1:sec:exercises} are where you
write the code and where the Cross-check shows you what three legs agreeing looks like, before a
later chapter shows you what it looks like when they do not.

% Section 1.1, from lectures/L01/appendix/a_circuits_and_units.md.

\section{Circuits, units, and loading}
\label{c1:app:a}

The circuit theory the rest of the book needs. With a background in electronics, read \secref{c1:sec:a1} for the conventions this book fixes, then skip to \secref{c1:sec:a6}.

\subsection{The sign convention, and why it is worth five minutes}
\label{c1:sec:a1}
Neither current direction nor voltage polarity is discoverable from a schematic. Fixed once:

\begin{itemize}
  \item \textbf{Current into a terminal is positive.}
  \item \textbf{A voltage is a difference.} $V_{AB}$ is the potential at A minus that at B; a single subscript such as $V_C$ is measured against ground.
  \item \textbf{Ground is node zero,} by definition rather than by measurement.
\end{itemize}
\textbf{A bias voltage that comes out negative, or an inverting stage whose gain comes out positive, is a convention error before it is a physics error.}

\subsection{The three quantities, and the one that is usually the answer}
\label{c1:sec:a2}
\begin{booktable}{@{}lllL@{}}
  \thead{Quantity} & \thead{Symbol} & \thead{Unit} & \thead{What it is} \\
  \midrule
  Current & $I$ & ampere & Charge per second past a point. \\
  Voltage & $V$ & volt & Energy per unit charge between two points. \\
  Resistance & $R$ & ohm & The ratio of the two, when that ratio is constant. \\
  Power & $P$ & watt & $VI$, and therefore $I^2R$ or $V^2/R$. \\
\end{booktable}

Ohm's law, $V = IR$, defines a resistor rather than stating a law of nature. Every device in Part~2 has a curve instead, and the work is choosing where on it to sit.

\textbf{Power decides whether a design survives a bench.} 10 V across 1 kilohm is 100 mW, which a common surface-mount part will not tolerate, and nothing in the schematic says so.

\subsection{Series, parallel, and the divider}
\label{c1:sec:a3}
Series adds resistance; parallel adds conductance.

\[
  R_{series} = R_1 + R_2
\]

\[
  \frac{1}{R_{parallel}} = \frac{1}{R_1} + \frac{1}{R_2}, \qquad R_{parallel} = \frac{R_1 R_2}{R_1 + R_2}
\]

\textbf{The parallel combination is always smaller than either part,} which catches most slips.

The \textbf{voltage divider} is the one circuit worth memorising: it is inside the bias network of Chapter~\ref{c6:ch:bias}, the feedback network of Chapter~\ref{c4:ch:feedback}, and the attenuator in front of every oscilloscope.

\[
  V_{out} = V_{in} \frac{R_{lower}}{R_{upper} + R_{lower}}
\]

\bookfigure{lectures/L01/appendix/images/divider.png}{A voltage divider from a 10 V supply, 33 kilohm over 6.8 kilohm, with the output node marked at 1.71 V and two annotations giving the Thevenin resistance of 5.6 kilohm and the loaded output of 1.09 V.}{c1:fig:divider}

33 kilohm over 6.8 kilohm on a 10 V supply:

\[
  V_{out} = 10 \times \frac{6800}{33000 + 6800} = 1.71\ \text{V}
\]

\textbf{1.71 V, not the 1.65 V it looks like it ought to be.} No E12 pair gives 1.65 V from 10 V here, and Chapter~\ref{c6:ch:bias} designs a transistor stage around the difference.

\subsection{Kirchhoff's two laws}
\label{c1:sec:a4}
The only physical content in circuit theory. \textbf{Current law:} currents into a node sum to zero, because charge does not accumulate at a junction.

\[
  \sum_k I_k = 0
\]

\textbf{Voltage law:} voltages around a closed loop sum to zero, because potential belongs to a point.

\[
  \sum_k V_k = 0
\]

\textbf{The current law is the one a program wants:} one equation per node, and a netlist has nodes.

\subsection{Thevenin and Norton}
\label{c1:sec:a5}
Any network of linear elements, seen from two terminals, behaves exactly like a voltage source in series with a resistance.

\begin{itemize}
  \item $V_{th}$ is the open-circuit voltage: what appears with nothing attached.
  \item $R_{th}$ is the resistance looking back in with every independent source killed. A voltage source becomes a short, a current source an open circuit.
\end{itemize}
Killing the divider's supply grounds the top of the upper leg, so the legs are in parallel:

\[
  R_{th} = R_{upper} \parallel R_{lower} = 5.64\ \text{k}\Omega
\]

\textbf{$R_{th}$ does not depend on which leg is the upper one, and is smaller than either.} Two megohms give one megohm; two ten-ohm resistors give five ohms.

\textbf{Norton} is the same source drawn as a current source in parallel with the same resistance, $I_{no} = V_{th}/R_{th}$, and is the form \secref{c1:app:b} wants.

\subsection{Loading, and why it decides everything later}
\label{c1:sec:a6}
The divider formula assumes nothing is attached. Attach something, and the lower leg becomes the lower leg in parallel with it. A 10 kilohm load on the 33k/6.8k divider:

\[
  V_{out} = 10 \times \frac{6800 \parallel 10000}{33000 + (6800 \parallel 10000)} = 1.09\ \text{V}
\]

\textbf{A third of the output gone, to a load most people would call light.} $R_{th}$ predicts it: 10 kilohm is not large compared with 5.6 kilohm.

\bookfigure{lectures/L01/appendix/images/divider_loading.png}{Divider output against load resistance, falling from 1.71 V unloaded to below 0.2 V at a 100 ohm load, with a marker showing that a load equal to the 5.6 kilohm Thevenin resistance halves the output.}{c1:fig:dividerloading}

\[
  V_{delivered} = V_{th} \frac{R_{load}}{R_{th} + R_{load}}
\]

A load equal to $R_{th}$ halves the output; ten times larger costs about 9 per cent; ten times smaller leaves about 9 per cent.

\textbf{This is the idea the book is built on.} Chapter~\ref{c3:ch:filters}: cascaded filter sections miss the corner you designed, because the second loads the first. Chapter~\ref{c8:ch:follower}: an emitter follower exists only to break a loading chain. Chapter~\ref{c10:ch:amplifier}: an op-amp loses 11 dB of open-loop gain to it, more than any other term.

\subsection{What this section is blind to}
\label{c1:sec:a7}
\begin{itemize}
  \item \textbf{Everything is linear.} From Chapter~\ref{c4:ch:feedback} the solver needs Newton-Raphson.
  \item \textbf{Everything is at DC.} No capacitor, no inductor, no frequency. Chapter~\ref{c2:ch:reactance} adds them by making this same solver complex, not by starting again.
  \item \textbf{Nothing has a tolerance.} Two 5 per cent resistors give a ratio good to about 10 per cent. Ignored until Chapter~\ref{c9:ch:diffpair}, where the matching of two transistors decides a stage's whole performance.
\end{itemize}

% Section 1.2, from lectures/L01/appendix/b_nodal_analysis.md.

\section{Nodal analysis, and what to build}
\label{c1:app:b}

The method that turns a circuit into a matrix, and the specification of the two components this chapter asks you to write. \secref{c1:sec:b5} is the specification; everything before it is what you need in order to write it.

\subsection{One unknown per node}
\label{c1:sec:b1}
Solving by hand means picking loops and current directions until the equations match the unknowns. The choices are yours rather than the circuit's, so it does not generalise into code.

\textbf{Nodal analysis makes them for you:} node voltages are the unknowns, the current law at each node supplies the equations, and the count comes out right. Ground is not one of them, so a circuit with $n$ nodes has $n - 1$ unknowns. \textbf{Five nodes is four unknowns.}

\subsection{The conductance stamp}
\label{c1:sec:b2}
For a resistor between nodes $a$ and $b$, the current leaving $a$ through it is

\[
  I_{a \to b} = \frac{V_a - V_b}{R} = G(V_a - V_b), \qquad G = \frac{1}{R}
\]

In $\mathbf{G}\mathbf{v} = \mathbf{i}$ one resistor contributes exactly four entries, always these four:

\[
  G_{aa} \mathrel{+}= G, \quad G_{bb} \mathrel{+}= G, \quad G_{ab} \mathrel{-}= G, \quad G_{ba} \mathrel{-}= G
\]

That pattern is a \textbf{stamp}, and it is the whole technique: never derive an equation, walk the element list, add each stamp, solve. Two consequences, both tested by the shipped suite:

\begin{itemize}
  \item \textbf{The diagonal accumulates.} Two resistors between the same pair of nodes stamp twice, giving the parallel combination without anyone computing one. A solver that assigns one matrix row per \emph{element} rather than per \emph{node} gets this wrong, and gets it wrong quietly.
  \item \textbf{Ground is skipped.} An entry whose row or column is ground is not written, which is why a resistor to ground contributes one entry rather than four.
\end{itemize}
A current source does not enter $\mathbf{G}$ at all; it adds to the right-hand side at its two nodes.

\subsection{The trouble with a voltage source}
\label{c1:sec:b3}
A voltage source has no conductance: its current is whatever the circuit demands, so it cannot be stamped. Add that current as an extra unknown and its constraint as an extra equation:

\[
  V_p - V_n = V_{source}
\]

One extra row and column per source, and the matrix stops being purely a conductance matrix. That is the whole of the ``modified'' in \textbf{modified nodal analysis}. The extra unknown says how much a supply delivers, which is how Chapter~\ref{c6:ch:bias} checks a bias point and Chapter~\ref{c10:ch:amplifier} a power budget.

\subsection{Two conventions this book fixes}
\label{c1:sec:b4}
Pinned by this chapter's test suite, because every component in the rest of the book reads them.

\textbf{A current source injects at its second node.} \code{addCurrentSource(from, to, current)} takes current out of \code{from} and pushes it into \code{to}. One milliamp into a node with one kilohm to ground puts that node at $+1$ V.

\textbf{A voltage source reports the current leaving its positive terminal.} \code{sourceCurrents[k]} flows out of source $k$'s positive terminal into the circuit; a 10 V source driving one kilohm reports $+10$ mA.

\textbf{That is the opposite sign to the raw modified-nodal unknown,} which comes out as the current flowing from the node into the source. Negating it is your job, and the suite checks that you did.

\subsection{What to build}
\label{c1:sec:b5}
Two headers. The paths are part of the specification: the test suite includes exactly these, and a component written anywhere else stays dormant however correct it is.

\subsubsection{\code{ael/net/netlist.hpp}}
A container. It holds elements and hands them to the solver, and it does nothing else.

\begin{booktable}{@{}LL@{}}
  \thead{Member} & \thead{Contract} \\
  \midrule
  \code{using Node = std::size_t} & A node index. \\
  \code{constexpr Node Ground{0}} & Node zero is ground, always. \\
  \code{addResistor(Node a, Node b, double resistance)} & Adds a resistor. \\
  \code{addCurrentSource(Node from, Node to, double current)} & Injects at \code{to}. See \secref{c1:sec:b4}. \\
  \code{addVoltageSource(Node positive, Node negative, double voltage)} & Fixes the difference. \\
  \code{nodeCount()} & Highest node index mentioned, plus one. An empty netlist returns 1. \\
  \code{resistorCount()}, \code{voltageSourceCount()}, \code{currentSourceCount()} & Element counts, by kind. \\
\end{booktable}

Nodes are mentioned rather than declared: there is no ``add a node'' call, because that would put a bookkeeping step in front of every exercise in the book.

The solver needs to read the elements back out. How you expose them is yours to decide; the tests only require that \code{solve} can be handed a \code{const Netlist&}.

\subsubsection{\code{ael/mna/solver.hpp}}
\begin{cppcode}
namespace ael::mna
{
struct Solution
{
    std::vector<double> nodeVoltages{};    ///< Indexed by Node. [Ground] is always 0.
    std::vector<double> sourceCurrents{};  ///< One per voltage source, in insertion order.
    bool solved{false};                    ///< False if the network has no unique solution.
};

[[nodiscard]] Solution solve(const net::Netlist& netlist) noexcept;
}
\end{cppcode}

Assemble from the stamps in \secref{c1:sec:b2} and \secref{c1:sec:b3}, then solve by Gaussian elimination with partial pivoting.

\textbf{\code{solved} is not decoration.} A floating node, or two voltage sources contradicting each other, gives a singular matrix. A solver that returns zeros there will be believed, and the reader spends an hour looking for the fault in their circuit rather than in their assumptions. Detect the singular pivot and say so. The suite has a test for exactly this.

\subsubsection{What good looks like}
Around 120 lines for both, of which the elimination is about 20. If yours is much longer, you are probably special-casing element kinds somewhere that a stamp would have handled.

\subsection{Why write a solver at all}
\label{c1:sec:b6}
\textbf{From Chapter~\ref{c7:ch:smallsignal} this book computes every stage twice:} once from a closed form you can reason about, once from a solver that knows nothing about amplifiers. Where they agree, the closed form is trustworthy and simpler. Where they disagree, one is dropping a term, and finding out which is the exercise. That only works if the solver is yours; a number from a tool you did not write is an authority claim.

\subsection{What this section is blind to}
\label{c1:sec:b7}
\begin{itemize}
  \item \textbf{Conditioning.} A direct solve handles the three or four orders of magnitude here without complaint. Chapter~\ref{c4:ch:feedback} brings a diode and twelve orders, at which point pivoting stops being a formality.
  \item \textbf{Sparsity.} Real simulators store the matrix sparsely and order it to limit fill-in. Nothing in this book is large enough for that to matter.
  \item \textbf{Anything time-varying.} This is a DC solve; Chapter~\ref{c2:ch:reactance} makes it complex, which covers one frequency in steady state. There is no transient analysis, so an underdamped filter's ringing is something Chapter~\ref{c3:ch:filters} computes rather than watches.
\end{itemize}

% Section 1.3, from lectures/L01/README.md: the objectives, the questions,
% the further reading, and what the next chapter does with all of it.


\section{Review}
\label{c1:sec:review}

\needspace{6\baselineskip}
\subsection*{What you should be able to do now}
After this chapter, you should be able to:
\begin{itemize}
  \item State the sign convention this book uses, and say what breaks if it is inverted.
  \item Compute a divider's output, its Thevenin resistance, and its output under any load, by hand.
  \item Say why a 10 kilohm load on a 33k/6.8k divider is not a light load.
  \item Write the node equation for any node in a resistive network without deriving it from scratch.
  \item Stamp a resistor, a current source and a voltage source into a matrix.
  \item Explain why a voltage source needs an extra unknown, and what that unknown is.
  \item Implement \code{ael::net} and \code{ael::mna} against the shipped suite.
\end{itemize}

\needspace{6\baselineskip}
\subsection*{Questions to test yourself}
You should be able to answer:
\begin{enumerate}
  \item A divider made of 1 megohm and 1 megohm, and one made of 10 ohms and 10 ohms, have the same ratio. What is different about them, and which would you put on an oscilloscope probe?
  \item Why is a divider's Thevenin resistance the same whichever leg you call the upper one?
  \item Nodal analysis gives one equation per node. A network has five nodes. How many unknowns are there, and why is it not five?
  \item You stamp a resistor between two nodes and get four matrix entries. Two are positive and two are negative. Which are which, and what would swapping them mean physically?
  \item A voltage source between two nodes cannot be written as a conductance. Why not, and what does the extra row in the matrix say?
  \item Your solver returns node voltages that are all zero for a network you know is live. Name two causes, and say which one the shipped suite catches.
  \item Superposition says two sources add. Your solver never mentions superposition. Why does it hold anyway, and in which chapter does it stop holding?
\end{enumerate}

\needspace{6\baselineskip}
\subsection*{Further reading}
\begin{itemize}
  \item \emph{The Art of Electronics}, Horowitz and Hill, chapter 1, for a second treatment of the same material at roughly the same level.
\end{itemize}

\needspace{6\baselineskip}
\subsection*{Looking ahead}
\begin{itemize}
  \item Chapter~\ref{c2:ch:reactance} makes the solver complex, which is the whole of the frequency domain and costs about thirty lines. Reactance, phasors, and the Bode plot as two straight lines and a corner.
\end{itemize}

% Section 1.4, from lectures/L01/appendix/c_exercises.md. The worked solutions are not
% printed; the aside says where they are.

\section{Exercises}
\label{c1:sec:exercises}

Eight, ending with the Cross-check. Do them in order; the last one needs the code from the two before it.

\begin{aside}
\textbf{How to check your work.} A \kindtag{Code} exercise is judged by the suite that ships with it: run \code{make test} at the root of the companion repository, with \code{AEL_DIR} pointing at your toolkit. A suite you have not started is reported as \code{SKIP} rather than as a failure, so the command is worth running from the first day. The hand calculations and designs are checked against the toolkit functions each one names.

\textbf{The solutions are not in this book.} They are published in full in the companion repository, in \file{lectures/L01/appendix}, including the plausible wrong answer where there is one. Read them once you have your own answers, including the ones you are unsure about, because being unsure and right is a different thing from being unsure and wrong and the solutions say which is which.
\end{aside}

% ------------------------------------------------------------------
\recall{The convention}
\label{c1:ex:1}
State the sign convention this book fixes for a current source and for a voltage source's reported current. Then answer this: your solver comes back with a divider output of $-1.71$ V instead of $+1.71$ V, and every voltage in the circuit is negated. Which of the two conventions did you get backwards, and how do you know it was that one rather than the other?

% ------------------------------------------------------------------
\recall{Two dividers}
\label{c1:ex:2}
A divider made of two 1 megohm resistors and a divider made of two 10 ohm resistors have the same ratio.

\begin{enumerate}
  \item What is the Thevenin resistance of each?
  \item One of them is the front end of an oscilloscope probe and the other is inside a power supply. Which is which, and why?
  \item A third divider uses two 10 ohm resistors across a 10 V supply. What is its power dissipation, and is that a problem?
\end{enumerate}

% ------------------------------------------------------------------
\handcalc{The divider, three ways}
\label{c1:ex:3}
For a 33 kilohm over 6.8 kilohm divider on a 10 V supply, by hand:

\begin{enumerate}
  \item The unloaded output.
  \item The Thevenin resistance.
  \item The output with a 10 kilohm load, computed twice: once by substituting the parallel combination into the divider formula, and once from the Thevenin equivalent. Confirm they agree, and say why they must.
  \item The load that would halve the output.
\end{enumerate}
\checkyourself{\code{divider}, \code{divider_output_resistance}}

% ------------------------------------------------------------------
\handcalc{Node equations}
\label{c1:ex:4}
A 1 mA current source injects into node 1. A 1 kilohm resistor runs from node 1 to ground, a 2.2 kilohm resistor from node 1 to node 2, and a 4.7 kilohm resistor from node 2 to ground.

\begin{enumerate}
  \item Write the current-law equation at node 1 and at node 2. Do not solve them yet.
  \item Write the two equations as a matrix, and identify each of the four entries as a stamp contribution from a named resistor.
  \item Solve for both node voltages.
  \item Kill the source and find the resistance looking into node 2. Now divide node 2's voltage from part 3 by the 1 mA source current. Those two numbers are both resistances and both about node 2. Say whether they agree, and if not, say precisely what each one is a resistance \emph{between}.
\end{enumerate}

% ------------------------------------------------------------------
\design{A reference divider}
\label{c1:ex:5}
Design a divider giving 2.5 V from a 10 V supply, subject to:

\begin{itemize}
  \item Thevenin resistance no greater than 2.2 kilohm, so the next stage does not load it.
  \item Divider current no greater than 1 mA, because the supply is a battery.
  \item E12 values only.
\end{itemize}
Give both resistors, the output voltage you actually achieve, the Thevenin resistance and the current. Then state the percentage error in the output, and say whether you would accept it.

There is a tempting wrong answer to this one. Find the pair whose \emph{ratio} is closest to correct, check it against both constraints, and say what it violates.

\checkyourself{\code{divider}, \code{divider_output_resistance}, \code{nearest_e12}}

% ------------------------------------------------------------------
\codeexercise{The netlist}
\label{c1:ex:6}
Implement \code{ael::net::Netlist} to the specification in \secref{c1:sec:b5}.

Four of this chapter's seventeen tests are for this. Getting them green needs no solver, so do this first and confirm the suite goes from four passing to eight.

% ------------------------------------------------------------------
\codeexercise{The solver}
\label{c1:ex:7}
Implement \code{ael::mna::solve} to the same specification.

Assemble the stamps, solve by Gaussian elimination with partial pivoting, and detect a singular matrix rather than returning zeros. All seventeen tests should pass.

Two of them are worth reading before you start rather than after you fail them: \code{Mna.ParallelResistorsCombine}, which fails for a solver that indexes by element instead of by node, and \code{Mna.FloatingNodeIsReportedRatherThanGuessed}, which fails for a solver that trusts its own elimination.

% ------------------------------------------------------------------
\crosscheck{The loaded divider}
\label{c1:ex:8}
The signature exercise of this book, and the shape every later one takes. Compute the same number three ways and reconcile them.

Take the 33 kilohm over 6.8 kilohm divider on a 10 V supply, loaded with 10 kilohm.

\begin{enumerate}
  \item \textbf{By hand.} Substitute the parallel combination into the divider formula. Write down the answer to three significant figures before doing anything else.
  \item \textbf{By your closed form.} Whatever you wrote for Exercise~\ref{c1:ex:3}, as a function, evaluated on the same numbers.
  \item \textbf{By your solver.} Build the netlist with four elements, one voltage source and three resistors, and read node 1 out of \code{ael::mna::solve}.
\end{enumerate}
Then reconcile. State the difference between each pair, and say what would have to be true for each difference to be non-zero.

\task{What to expect}
\textbf{Legs 1 and 2 should agree to every digit you wrote down.} They are the same formula evaluated by two different pieces of hardware. A disagreement here is an arithmetic slip in leg 1 or a typo in leg 2, and it is not interesting except that finding it is faster than finding it later.

\textbf{Legs 2 and 3 should agree to about twelve significant figures.} Both are solving the same linear problem in double precision, one by algebra you did in advance and one by elimination. The difference is rounding and nothing else, and it should be somewhere around $10^{-13}$ relative.

\textbf{A difference of a few per cent means one of them is not solving the circuit you think it is.} The usual cause is the load stamped between the wrong pair of nodes, which quietly gives you a different circuit rather than an error. The second usual cause is the supply node and the output node swapped, which gives you the reciprocal divider.

\textbf{A difference of exactly a factor of two, or exactly the unloaded answer, means the load is not in the circuit at all.} Check \code{resistorCount()}.

This exercise is deliberately one where all three legs agree. That is worth doing once, so that you know what agreement looks like before Chapter~\ref{c7:ch:smallsignal}, where a closed form and the solver disagree by a factor of ten and the closed form is the one that is wrong.

